% Source: book pp. 18--26 (PDF 32--40); solutions pp. 14--24, 168.
\section{Subspaces}

By considering subspaces, we can greatly expand our examples of vector spaces.

\begin{definition}[Subspace]\label{ax:1.33}
A subset $U$ of $V$ is called a \emph{subspace} of $V$ if $U$ is also a vector
space with the same additive identity, addition, and scalar multiplication as
on $V$.
\end{definition}

\marginnote{Some people use the terminology \emph{linear subspace}, which means
the same as subspace.}%
The next result gives the easiest way to check whether a subset of a vector
space is a subspace.

\begin{result}[Conditions for a subspace]\label{ax:1.34}
A subset $U$ of $V$ is a subspace of $V$ if and only if $U$ satisfies the
following three conditions.
\begin{axprops}
\item[additive identity] $0\in U$.
\item[closed under addition] $u,w\in U$ implies $u+w\in U$.
\item[closed under scalar multiplication] $a\in\F$ and $u\in U$ implies
  $au\in U$.
\end{axprops}
\end{result}

\begin{proof}
If $U$ is a subspace of $V$, then $U$ satisfies the three conditions above by
the definition of vector space.

Conversely, suppose $U$ satisfies the three conditions above. The first
condition ensures that the additive identity of $V$ is in $U$. The second
condition ensures that addition makes sense on $U$. The third condition ensures
that scalar multiplication makes sense on $U$.

\marginnote{The additive identity condition above could be replaced with the
condition that $U$ is nonempty (because then taking $u\in U$ and multiplying it
by $0$ would imply that $0\in U$). However, if a subset $U$ of $V$ is indeed a
subspace, then usually the quickest way to show that $U$ is nonempty is to show
that $0\in U$.}%
If $u\in U$, then $-u$ [which equals $(-1)u$ by \ref{ax:1.32}] is also in $U$
by the third condition above. Hence every element of $U$ has an additive
inverse in $U$.

The other parts of the definition of a vector space, such as associativity and
commutativity, are automatically satisfied for $U$ because they hold on the
larger space $V$. Thus $U$ is a vector space and hence is a subspace of $V$.
\end{proof}

The three conditions in the result above usually enable us to determine quickly
whether a given subset of $V$ is a subspace of $V$. You should verify all
assertions in the next example.

\begin{example}[Subspaces]\label{ax:1.35}\leavevmode
\begin{parts}
\item If $b\in\F$, then
  \[ \{(x_1,x_2,x_3,x_4)\in\F^4 : x_3=5x_4+b\} \]
  is a subspace of $\F^4$ if and only if $b=0$.
\item The set of continuous real-valued functions on the interval $[0,1]$ is a
  subspace of $\R^{[0,1]}$.
\item The set of differentiable real-valued functions on $\R$ is a subspace of
  $\R^{\R}$.
\item The set of differentiable real-valued functions $f$ on the interval
  $(0,3)$ such that $f'(2)=b$ is a subspace of $\R^{(0,3)}$ if and only if
  $b=0$.
\item The set of all sequences of complex numbers with limit $0$ is a subspace
  of $\C^\infty$.
\end{parts}
\end{example}

\marginnote{The set $\{0\}$ is the smallest subspace of $V$, and $V$ itself is
the largest subspace of $V$. The empty set is not a subspace of $V$ because a
subspace must be a vector space and hence must contain at least one element,
namely, an additive identity.}%
Verifying some of the items above shows the linear structure underlying parts
of calculus. For example, (b) above requires the result that the sum of two
continuous functions is continuous. As another example, (d) above requires the
result that for a constant $c$, the derivative of $cf$ equals $c$ times the
derivative of $f$.

The subspaces of $\R^2$ are precisely $\{0\}$, all lines in $\R^2$ containing
the origin, and $\R^2$. The subspaces of $\R^3$ are precisely $\{0\}$, all lines
in $\R^3$ containing the origin, all planes in $\R^3$ containing the origin, and
$\R^3$. To prove that all these objects are indeed subspaces is
straightforward---the hard part is to show that they are the only subspaces of
$\R^2$ and $\R^3$. That task will be easier after we introduce some additional
tools in the next chapter.

\subsection{Sums of Subspaces}

\marginnote{The union of subspaces is rarely a subspace (see Exercise~12),
which is why we usually work with sums rather than unions.}%
When dealing with vector spaces, we are usually interested only in subspaces,
as opposed to arbitrary subsets. The notion of the sum of subspaces will be
useful.

\begin{definition}[Sum of subspaces]\label{ax:1.36}
Suppose $V_1,\dots,V_m$ are subspaces of $V$. The \emph{sum} of
$V_1,\dots,V_m$, denoted by $V_1+\cdots+V_m$, is the set of all possible sums
of elements of $V_1,\dots,V_m$. More precisely,
\[ V_1+\cdots+V_m = \{v_1+\cdots+v_m : v_1\in V_1,\dots,v_m\in V_m\}. \]
\end{definition}

Let's look at some examples of sums of subspaces.

\begin{example}[A sum of subspaces of $\F^3$]\label{ax:1.37}
Suppose $U$ is the set of all elements of $\F^3$ whose second and third
coordinates equal $0$, and $W$ is the set of all elements of $\F^3$ whose first
and third coordinates equal $0$:
\[ U = \{(x,0,0)\in\F^3 : x\in\F\} \quad\text{and}\quad
   W = \{(0,y,0)\in\F^3 : y\in\F\}. \]
Then
\[ U+W = \{(x,y,0)\in\F^3 : x,y\in\F\}, \]
as you should verify.
\end{example}

\begin{example}[A sum of subspaces of $\F^4$]\label{ax:1.38}
Suppose
\[ U = \{(x,x,y,y)\in\F^4 : x,y\in\F\} \quad\text{and}\quad
   W = \{(x,x,x,y)\in\F^4 : x,y\in\F\}. \]
Using words rather than symbols, we could say that $U$ is the set of elements
of $\F^4$ whose first two coordinates equal each other and whose third and
fourth coordinates equal each other. Similarly, $W$ is the set of elements of
$\F^4$ whose first three coordinates equal each other.

To find a description of $U+W$, consider a typical element $(a,a,b,b)$ of $U$
and a typical element $(c,c,c,d)$ of $W$, where $a,b,c,d\in\F$. We have
\[ (a,a,b,b)+(c,c,c,d) = (a+c,a+c,b+c,b+d), \]
which shows that every element of $U+W$ has its first two coordinates equal to
each other. Thus\refstepcounter{theorem}
\[ U+W \subseteq \{(x,x,y,z)\in\F^4 : x,y,z\in\F\}.
   \tag*{\thetheorem}\label{ax:1.39} \]

To prove the inclusion in the other direction, suppose $x,y,z\in\F$. Then
\[ (x,x,y,z) = (x,x,y,y)+(0,0,0,z-y), \]
where the first vector on the right is in $U$ and the second vector on the
right is in $W$. Thus $(x,x,y,z)\in U+W$, showing that the inclusion
\ref{ax:1.39} also holds in the opposite direction. Hence
\[ U+W = \{(x,x,y,z)\in\F^4 : x,y,z\in\F\}, \]
which shows that $U+W$ is the set of elements of $\F^4$ whose first two
coordinates equal each other.
\end{example}

The next result states that the sum of subspaces is a subspace, and is in fact
the smallest subspace containing all the summands (which means that every
subspace containing all the summands also contains the sum).

\begin{result}[Sum of subspaces is the smallest containing subspace]\label{ax:1.40}
Suppose $V_1,\dots,V_m$ are subspaces of $V$. Then $V_1+\cdots+V_m$ is the
smallest subspace of $V$ containing $V_1,\dots,V_m$.
\end{result}

\begin{proof}
The reader can verify that $V_1+\cdots+V_m$ contains the additive identity $0$
and is closed under addition and scalar multiplication. Thus \ref{ax:1.34}
implies that $V_1+\cdots+V_m$ is a subspace of $V$.

\marginnote{Sums of subspaces in the theory of vector spaces are analogous to
unions of subsets in set theory. Given two subspaces of a vector space, the
smallest subspace containing them is their sum. Analogously, given two subsets
of a set, the smallest subset containing them is their union.}%
The subspaces $V_1,\dots,V_m$ are all contained in $V_1+\cdots+V_m$ (to see
this, consider sums $v_1+\cdots+v_m$ where all except one of the $v_k$'s are
$0$). Conversely, every subspace of $V$ containing $V_1,\dots,V_m$ contains
$V_1+\cdots+V_m$ (because subspaces must contain all finite sums of their
elements). Thus $V_1+\cdots+V_m$ is the smallest subspace of $V$ containing
$V_1,\dots,V_m$.
\end{proof}

\subsection{Direct Sums}

Suppose $V_1,\dots,V_m$ are subspaces of $V$. Every element of
$V_1+\cdots+V_m$ can be written in the form
\[ v_1+\cdots+v_m, \]
where each $v_k\in V_k$. Of special interest are cases in which each vector in
$V_1+\cdots+V_m$ can be represented in the form above in only one way. This
situation is so important that it gets a special name (direct sum) and a
special symbol ($\oplus$).

\begin{definition}[Direct sum, $\oplus$]\label{ax:1.41}
Suppose $V_1,\dots,V_m$ are subspaces of $V$.
\begin{itemize}
\item The sum $V_1+\cdots+V_m$ is called a \emph{direct sum} if each element of
  $V_1+\cdots+V_m$ can be written in only one way as a sum $v_1+\cdots+v_m$,
  where each $v_k\in V_k$.
\item If $V_1+\cdots+V_m$ is a direct sum, then $V_1\oplus\cdots\oplus V_m$
  denotes $V_1+\cdots+V_m$, with the $\oplus$ notation serving as an indication
  that this is a direct sum.
\end{itemize}
\end{definition}

\begin{example}[A direct sum of two subspaces]\label{ax:1.42}
Suppose $U$ is the subspace of $\F^3$ of those vectors whose last coordinate
equals $0$, and $W$ is the subspace of $\F^3$ of those vectors whose first two
coordinates equal $0$:
\[ U = \{(x,y,0)\in\F^3 : x,y\in\F\} \quad\text{and}\quad
   W = \{(0,0,z)\in\F^3 : z\in\F\}. \]
Then $\F^3=U\oplus W$, as you should verify.
\end{example}

\marginnote[6mm]{To produce $\oplus$ in \TeX, type \texttt{\string\oplus}.}%
\begin{example}[A direct sum of multiple subspaces]\label{ax:1.43}
Suppose $V_k$ is the subspace of $\F^n$ of those vectors whose coordinates are
all $0$, except possibly in the $k^{\text{th}}$ slot; for example,
$V_2=\{(0,x,0,\dots,0)\in\F^n : x\in\F\}$. Then
\[ \F^n = V_1\oplus\cdots\oplus V_n, \]
as you should verify.
\end{example}

Sometimes nonexamples add to our understanding as much as examples.

\begin{example}[A sum that is not a direct sum]\label{ax:1.44}
Suppose
\begin{align*}
V_1 &= \{(x,y,0)\in\F^3 : x,y\in\F\},\\
V_2 &= \{(0,0,z)\in\F^3 : z\in\F\},\\
V_3 &= \{(0,y,y)\in\F^3 : y\in\F\}.
\end{align*}
Then $\F^3=V_1+V_2+V_3$ because every vector $(x,y,z)\in\F^3$ can be written as
\[ (x,y,z) = (x,y,0)+(0,0,z)+(0,0,0), \]
where the first vector on the right side is in $V_1$, the second vector is in
$V_2$, and the third vector is in $V_3$.

However, $\F^3$ does not equal the direct sum of $V_1,V_2,V_3$, because the
vector $(0,0,0)$ can be written in more than one way as a sum $v_1+v_2+v_3$,
with each $v_k\in V_k$. Specifically, we have
\[ (0,0,0) = (0,1,0)+(0,0,1)+(0,-1,-1) \]
and, of course,
\[ (0,0,0) = (0,0,0)+(0,0,0)+(0,0,0), \]
where the first vector on the right side of each equation above is in $V_1$,
the second vector is in $V_2$, and the third vector is in $V_3$. Thus the sum
$V_1+V_2+V_3$ is not a direct sum.
\end{example}

\marginnote[-8mm]{The symbol $\oplus$, which is a plus sign inside a circle, reminds
us that we are dealing with a special type of sum of subspaces---each element
in the direct sum can be represented in only one way as a sum of elements from
the specified subspaces.}%
The definition of direct sum requires every vector in the sum to have a unique
representation as an appropriate sum. The next result shows that when deciding
whether a sum of subspaces is a direct sum, we only need to consider whether
$0$ can be uniquely written as an appropriate sum.

\begin{result}[Condition for a direct sum]\label{ax:1.45}
Suppose $V_1,\dots,V_m$ are subspaces of $V$. Then $V_1+\cdots+V_m$ is a direct
sum if and only if the only way to write $0$ as a sum $v_1+\cdots+v_m$, where
each $v_k\in V_k$, is by taking each $v_k$ equal to $0$.
\end{result}

\begin{proof}
First suppose $V_1+\cdots+V_m$ is a direct sum. Then the definition of direct
sum implies that the only way to write $0$ as a sum $v_1+\cdots+v_m$, where each
$v_k\in V_k$, is by taking each $v_k$ equal to $0$.

Now suppose that the only way to write $0$ as a sum $v_1+\cdots+v_m$, where
each $v_k\in V_k$, is by taking each $v_k$ equal to $0$. To show that
$V_1+\cdots+V_m$ is a direct sum, let $v\in V_1+\cdots+V_m$. We can write
\[ v = v_1+\cdots+v_m \]
for some $v_1\in V_1,\dots,v_m\in V_m$. To show that this representation is
unique, suppose we also have
\[ v = u_1+\cdots+u_m, \]
where $u_1\in V_1,\dots,u_m\in V_m$. Subtracting these two equations, we have
\[ 0 = (v_1-u_1)+\cdots+(v_m-u_m). \]
Because $v_1-u_1\in V_1,\dots,v_m-u_m\in V_m$, the equation above implies that
each $v_k-u_k$ equals $0$. Thus $v_1=u_1,\dots,v_m=u_m$, as desired.
\end{proof}

\marginnote{The symbol $\iff$ used below means ``if and only if''; this symbol
could also be read to mean ``is equivalent to''.}%
The next result gives a simple condition for testing whether a sum of two
subspaces is a direct sum.

\begin{result}[Direct sum of two subspaces]\label{ax:1.46}
Suppose $U$ and $W$ are subspaces of $V$. Then
\[ U+W \text{ is a direct sum} \iff U\cap W=\{0\}. \]
\end{result}

\begin{proof}
First suppose that $U+W$ is a direct sum. If $v\in U\cap W$, then
$0=v+(-v)$, where $v\in U$ and $-v\in W$. By the unique representation of $0$
as the sum of a vector in $U$ and a vector in $W$, we have $v=0$. Thus
$U\cap W=\{0\}$, completing the proof in one direction.

To prove the other direction, now suppose $U\cap W=\{0\}$. To prove that $U+W$
is a direct sum, suppose $u\in U$, $w\in W$, and
\[ 0 = u+w. \]
To complete the proof, we only need to show that $u=w=0$ (by \ref{ax:1.45}).
The equation above implies that $u=-w\in W$. Thus $u\in U\cap W$. Hence $u=0$,
which by the equation above implies that $w=0$, completing the proof.
\end{proof}

\marginnote{Sums of subspaces are analogous to unions of subsets. Similarly,
direct sums of subspaces are analogous to disjoint unions of subsets. No two
subspaces of a vector space can be disjoint, because both contain $0$. So
disjointness is replaced, at least in the case of two subspaces, with the
requirement that the intersection equal $\{0\}$.}%
The result above deals only with the case of two subspaces. When asking about
a possible direct sum with more than two subspaces, it is not enough to test
that each pair of the subspaces intersect only at $0$. To see this, consider
Example~\ref{ax:1.44}. In that nonexample of a direct sum, we have
$V_1\cap V_2=V_1\cap V_3=V_2\cap V_3=\{0\}$.

% ------------------------------------------------------------------ exercises
\exercisesheading{1C}

\begin{exercise}
For each of the following subsets of $\F^3$, determine whether it is a subspace
of $\F^3$.
\begin{parts}
\item $\{(x_1,x_2,x_3)\in\F^3 : x_1+2x_2+3x_3=0\}$
\item $\{(x_1,x_2,x_3)\in\F^3 : x_1+2x_2+3x_3=4\}$
\item $\{(x_1,x_2,x_3)\in\F^3 : x_1x_2x_3=0\}$
\item $\{(x_1,x_2,x_3)\in\F^3 : x_1=5x_3\}$
\end{parts}
\end{exercise}
\begin{solution}
Let $U$ be the subset of $\F^3$ in question.
\begin{parts}
\item Since the zero vector $(0,0,0)$ satisfies $0+2(0)+3(0)=0$, $U$ is
  nonempty. Let $u=(u_1,u_2,u_3)$ and $v=(v_1,v_2,v_3)$ be elements of $U$.
  Then
  \[ (u_1+v_1)+2(u_2+v_2)+3(u_3+v_3)
     = (u_1+2u_2+3u_3)+(v_1+2v_2+3v_3) = 0+0 = 0. \]
  Hence, $u+v$ is in $U$. Let $a\in\F$. Then
  \[ au_1+2(au_2)+3(au_3) = a(u_1+2u_2+3u_3) = a\cdot0 = 0. \]
  Hence, $au$ is in $U$. Therefore, $U$ is a subspace of $\F^3$.
\item The zero vector $(0,0,0)$ does not satisfy $0+2(0)+3(0)=4$, so $U$ does
  not contain the zero vector. Hence, it is not a subspace of $\F^3$.
\item Suppose $u=(1,0,0)$ and $v=(0,1,1)$. Then $u$ and $v$ are in $U$ since
  $1\cdot0\cdot0=0$ and $0\cdot1\cdot1=0$. However, $u+v=(1,1,1)$ is not in $U$
  since $1\cdot1\cdot1=1\ne0$. Hence, $U$ is not closed under addition and is
  not a subspace of $\F^3$.
\item Since the zero vector $(0,0,0)$ satisfies $0=5(0)$, $U$ is nonempty. Let
  $u=(u_1,u_2,u_3)$ and $v=(v_1,v_2,v_3)$ be elements of $U$. Then
  \[ u_1+v_1 = 5u_3+5v_3 = 5(u_3+v_3). \]
  Hence, $u+v$ is in $U$. Let $a\in\F$. Then
  \[ au_1 = a(5u_3) = 5(au_3). \]
  Hence, $au$ is in $U$. Therefore, $U$ is a subspace of $\F^3$.
\end{parts}
\end{solution}

\begin{exercise}
Verify all assertions about subspaces in Example~\ref{ax:1.35}.
\end{exercise}
\begin{solution}
Let $U$ be the subset of $\F^4$ in question.
\begin{parts}
\item Suppose $U$ is a subspace of $\F^4$. Then it must contain the zero vector
  $(0,0,0,0)$. Then $0=5(0)+b$, which implies $b=0$.

  Conversely, suppose $b=0$. Then
  $U=\{(x_1,x_2,x_3,x_4)\in\F^4 \mid x_3=5x_4\}$. The proof is similar to that
  of Exercise~1C1~(d), hence $U$ is a subspace of $\F^4$.
\item Note that the $0$ function on $[0,1]$ is continuous, so $0\in U$. Let
  $f,g\in U$ be continuous functions. Then $f+g$ is continuous, so $f+g\in U$.
  Let $a\in\R$. Then $af$ is continuous, so $af\in U$. Therefore, $U$ is a
  subspace of $\R^{[0,1]}$.
\item Since the derivative of the $0$ function is the $0$ function, $0\in U$.
  Let $f,g\in U$ be differentiable functions. Then $f+g$ is differentiable, so
  $f+g\in U$. Let $a\in\R$. Then $af$ is differentiable, so $af\in U$.
  Therefore, $U$ is a subspace of $\R^{\R}$.
\item Suppose $U$ is a subspace of $\R^{(0,3)}$. Then it must contain the $0$
  function on $(0,3)$. Then $0'=0=b$, which implies $b=0$.

  Conversely, suppose $b=0$. Then $U=\{f\in\R^{(0,3)} \mid f'(2)=0\}$. Let
  $f,g\in U$ be differentiable functions such that $f'(2)=g'(2)=0$. Then
  \[ (f+g)'(2) = f'(2)+g'(2) = 0+0 = 0, \]
  so $f+g\in U$. Let $a\in\R$. Then
  \[ (af)'(2) = af'(2) = a\cdot0 = 0, \]
  so $af\in U$. Therefore, $U$ is a subspace of $\R^{(0,3)}$.
\item Let $U$ be the set of all sequences of complex numbers with limit $0$.
  Then the zero sequence is in $U$. Let $(z_n),(w_n)\in U$ be sequences with
  limit $0$. Then
  \[ \lim_{n\to\infty}(z_n+w_n) = \lim_{n\to\infty}z_n+\lim_{n\to\infty}w_n
     = 0+0 = 0, \]
  so $(z_n)+(w_n)\in U$. Let $a\in\C$. Then
  \[ \lim_{n\to\infty}(az_n) = a\lim_{n\to\infty}z_n = a\cdot0 = 0, \]
  so $a(z_n)\in U$. Therefore, $U$ is a subspace of $\C^\infty$.
\end{parts}
\end{solution}

\begin{exercise}
Show that the set of differentiable real-valued functions $f$ on the interval
$(-4,4)$ such that $f'(-1)=3f(2)$ is a subspace of $\R^{(-4,4)}$.
\end{exercise}
\begin{solution}
Let $U$ be the set in question. Then the zero function is in $U$ since
$0'(-1)=0=3\cdot0=3\cdot0(2)$. Let $f,g\in U$ be differentiable functions such
that $f'(-1)=3f(2)$ and $g'(-1)=3g(2)$. Then
\[ (f+g)'(-1) = f'(-1)+g'(-1) = 3f(2)+3g(2) = 3\bigl(f(2)+g(2)\bigr)
   = 3(f+g)(2), \]
so $f+g\in U$. Let $a\in\R$. Then
\[ (af)'(-1) = af'(-1) = a\cdot3f(2) = 3(af)(2), \]
so $af\in U$. Therefore, $U$ is a subspace of $\R^{(-4,4)}$.
\end{solution}

\begin{exercise}
Suppose $b\in\R$. Show that the set of continuous real-valued functions $f$ on
the interval $[0,1]$ such that $\int_0^1 f=b$ is a subspace of $\R^{[0,1]}$ if
and only if $b=0$.
\end{exercise}
\begin{solution}
Let $U$ be the set in question. Suppose $U$ is a subspace of $\R^{[0,1]}$. Then
it must contain the zero function on $[0,1]$. Then
\[ \int_0^1 0 = 0 = b, \]
which implies $b=0$.

Conversely, suppose $b=0$. Then
\[ U = \biggl\{f\in\R^{[0,1]} \biggm| \int_0^1 f=0\biggr\}. \]
The zero function is in $U$ since $\int_0^1 0=0$. Let $f,g\in U$ be continuous
functions such that $\int_0^1 f=\int_0^1 g=0$. Then
\[ \int_0^1 (f+g) = \int_0^1 f+\int_0^1 g = 0+0 = 0, \]
so $f+g\in U$. Let $a\in\R$. Then
\[ \int_0^1 (af) = a\int_0^1 f = a\cdot0 = 0, \]
so $af\in U$. Therefore, $U$ is a subspace of $\R^{[0,1]}$.
\end{solution}

\begin{exercise}
Is $\R^2$ a subspace of the complex vector space $\C^2$?
\end{exercise}
\begin{solution}
No, $\R^2$ is not a subspace of $\C^2$. Let $u=(1,0)\in\R^2$ and $a=i\in\C$.
Then $au=(i,0)\notin\R^2$, so $\R^2$ is not closed under scalar multiplication
and is not a subspace of $\C^2$.
\end{solution}

\begin{exercise}\leavevmode
\begin{parts}
\item Is $\{(a,b,c)\in\R^3 : a^3=b^3\}$ a subspace of $\R^3$?
\item Is $\{(a,b,c)\in\C^3 : a^3=b^3\}$ a subspace of $\C^3$?
\end{parts}
\end{exercise}
\begin{solution}
\solnote{Manual Remark}{Here, I present an elementary argument that does not
require the use of De Moivre's Theorem; rather, I arrive at the conclusion
directly. However, readers who are familiar with it may find it easier to note
that the complex number $1$ has three cube roots of unity. Other readers may
also use Vieta's formulas to arrive at the same polynomial in $\omega$ to obtain
$1+\omega=-\omega^2$.}
Let $U$ be the subset in question.
\begin{parts}
\item Since the zero vector $(0,0,0)$ satisfies $0^3=0^3$, $U$ is nonempty. Let
  $x=(x_1,x_2,x_3)$ and $y=(y_1,y_2,y_3)$ be elements of $U$. Then
  $x_1^3=x_2^3$ and $y_1^3=y_2^3$ both imply that $x_1=x_2$ and $y_1=y_2$.
  Hence,
  \[ (x_1+y_1)^3 = (x_2+y_2)^3, \]
  so $x+y\in U$. Let $a\in\R$. Then
  \[ (ax_1)^3 = a^3x_1^3 = a^3x_2^3 = (ax_2)^3, \]
  so $ax\in U$. Therefore, $U$ is a subspace of $\R^3$.
\item No, $U$ is not a subspace of $\C^3$. To see this, consider the cubic
  polynomial $x^3-1=0$. This polynomial has three distinct roots in $\C$ as
  follows:
  \[ x^3-1 = (x-1)(x^2+x+1) = 0 \]
  Applying the quadratic formula to $x^2+x+1=0$, we get
  \[ x = \frac{-1\pm\sqrt{(-1)^2-4(1)(1)}}{2(1)} = \frac{-1\pm i\sqrt3}{2}. \]
  Let
  \[ \omega = \frac{-1+i\sqrt3}{2}. \]
  Then $\omega^3=1$, and
  \[ \omega^2 = \frac{-1-i\sqrt3}{2}. \]
  Moreover, it is not hard to see that $\omega+\omega^2=-1$. Let $u=(1,1,0)$
  and $v=(\omega,1,0)$. Then $u,v\in U$ since $1^3=1^3$ and $\omega^3=1^3$.
  Then $u+v=(1+\omega,2,0)$. However,
  \[ (1+\omega)^3 = 1+3\omega+3\omega^2+\omega^3 = 1+3(-1)+1 = -1 \ne 2^3 = 8. \]
  Hence, $u+v\notin U$, so $U$ is not closed under addition and is not a
  subspace of $\C^3$.
\end{parts}
\end{solution}

\begin{exercise}
Prove or give a counterexample: If $U$ is a nonempty subset of $\R^2$ such that
$U$ is closed under addition and under taking additive inverses (meaning
$-u\in U$ whenever $u\in U$), then $U$ is a subspace of $\R^2$.
\end{exercise}
\begin{solution}
Consider the subset $\Q^2$ of $\R^2$. Then $\Q^2$ is nonempty because
$(0,0)\in\Q^2$, closed under addition since the sum of any two rational numbers
is rational, and closed under taking additive inverses since the additive
inverse of a rational number is rational. However, for any irrational number
$r\in\R$, the vector $(r,0)=r(1,0)\in\R^2$ is not in $\Q^2$. Hence, $\Q^2$ is
not closed under scalar multiplication and is not a subspace of $\R^2$.
Therefore, the statement is false.
\end{solution}

\begin{exercise}
Give an example of a nonempty subset $U$ of $\R^2$ such that $U$ is closed
under scalar multiplication, but $U$ is not a subspace of $\R^2$.
\end{exercise}
\begin{solution}
Consider the subset
\[ U = \{(x,0) \mid x\in\R\}\cup\{(0,y) \mid y\in\R\}. \]
Then $U$ is nonempty since $(1,0)\in U$, and it is closed under scalar
multiplication since any scalar multiple of $(1,0)$ or $(0,1)$ is still in
$U$. However, $U$ is not closed under addition since
$(1,0)+(0,1)=(1,1)\notin U$. Hence, $U$ is not a subspace of $\R^2$.
\end{solution}

\begin{exercise}
A function $f\colon\R\to\R$ is called \emph{periodic} if there exists a
positive number $p$ such that $f(x)=f(x+p)$ for all $x\in\R$. Is the set of
periodic functions from $\R$ to $\R$ a subspace of $\R^{\R}$? Explain.
\end{exercise}
\begin{solution}
It is not a subspace of $\R^{\R}$. Consider the periodic functions
$f(x)=\cos x$ and $g(x)=\cos\sqrt2\,x$. Then the sum $f+g$ attains a maximum
value of $2$ at $0$. If $f+g$ were periodic with period $p$, then
$(f+g)(p)=(f+g)(0)=2$. This implies $\cos p=\cos p\sqrt2=1$. Then $p=2\pi m$
and $p\sqrt2=2\pi n$ for positive integers $m$ and $n$. Since $p\sqrt2=2\pi n$
implies $p=\sqrt2\pi n$, we have
\[ 2\pi m = \sqrt2\pi n \implies \sqrt2 = \frac nm, \]
which is a contradiction, since $\sqrt2$ is irrational. Then $f+g$ is not
periodic.
\end{solution}

\begin{exercise}
Suppose $V_1$ and $V_2$ are subspaces of $V$. Prove that the intersection
$V_1\cap V_2$ is a subspace of $V$.
\end{exercise}
\begin{solution}
Let $U=V_1\cap V_2$. Since both $V_1$ and $V_2$ are subspaces of $V$, they both
contain the zero vector of $V$. Hence, $0\in U$, so $U$ is nonempty. Let
$u,v\in U$. Then $u,v\in V_1$ and $u,v\in V_2$. Since both $V_1$ and $V_2$ are
subspaces, they are closed under addition. Therefore, $u+v\in V_1$ and
$u+v\in V_2$, which implies that $u+v\in U$. Let $a\in\F$. Then since both
$V_1$ and $V_2$ are closed under scalar multiplication, we have that
$au\in V_1$ and $au\in V_2$, which implies that $au\in U$. Therefore, $U$ is a
subspace of $V$.
\end{solution}

\begin{exercise}
Prove that the intersection of every collection of subspaces of $V$ is a
subspace of $V$.
\end{exercise}
\begin{solution}
Let $\mathcal{C}$ be a collection of subspaces of $V$, and let
\[ U = \bigcap_{W\in\mathcal{C}} W. \]
Since each subspace in $\mathcal{C}$ contains the zero vector, $0\in U$, so $U$
is nonempty. Let $u,v\in U$. Then $u,v\in W$ for all $W\in\mathcal{C}$. Since
each $W$ is a subspace, they are closed under addition, so $u+v\in W$ for all
$W\in\mathcal{C}$, which implies $u+v\in U$. Let $a\in\F$. Then $au\in W$ for
all $W\in\mathcal{C}$, which implies $au\in U$. Therefore, $U$ is a subspace of
$V$.
\end{solution}

\begin{exercise}
Prove that the union of two subspaces of $V$ is a subspace of $V$ if and only
if one of the subspaces is contained in the other.
\end{exercise}
\begin{solution}
Let $U,W$ be subspaces of $V$. Suppose $U\subseteq W$. Then $U\cup W=W$, which
is a subspace of $V$. Similarly, if $W\subseteq U$, then $U\cup W=U$, which is
a subspace of $V$.

Conversely, suppose $U\cup W$ is a subspace of $V$. If neither $U\subseteq W$
nor $W\subseteq U$, then there exist $u\in U\setminus W$ and
$w\in W\setminus U$. Since $U\cup W$ is a subspace, it must be closed under
addition. Then $u+w\in U\cup W$. However, $u+w\notin U$ because if it were,
then $w=(u+w)-u\in U$, which is a contradiction. Similarly, $u+w\notin W$
because if it were, then $u=(u+w)-w\in W$, which is also a contradiction.
Hence, $u+w\notin U\cup W$, which contradicts the assumption that $U\cup W$ is
a subspace. Therefore, one of the subspaces must be contained in the other.
\end{solution}

\begin{exercise}
Prove that the union of three subspaces of $V$ is a subspace of $V$ if and only
if one of the subspaces contains the other two.
\exnote{This exercise is surprisingly harder than Exercise~12, possibly because
this exercise is not true if we replace $\F$ with a field containing only two
elements.}
\end{exercise}
\begin{solution}
\solnote{Axler Note}{This result fails when the field has only two elements.}
\solnote{Manual Remark}{The proof of this exercise is more involved than the
proof of Exercise~1C12 in that it requires a more careful analysis of the
possible relationships between the three subspaces.}
\solnote{Appendix}{See the alternate proof at the end of this solution for a
more elegant proof of this result, which implicitly uses the result that
$\operatorname{char}(\F)\ne2$.}
Let $U_1,U_2,U_3$ be subspaces of $V$. Without loss of generality, suppose
$U_2\subseteq U_1$ and $U_3\subseteq U_1$. Then $U_1\cup U_2\cup U_3=U_1$,
which is a subspace of $V$.

Conversely, suppose $U_1\cup U_2\cup U_3$ is a subspace of $V$. Suppose that
some pair satisfies containment, say $U_3\subseteq U_2$. Let
$W=U_2\cup U_3=U_2$. Then $U_1\cup W=U_1\cup U_2\cup U_3$. By Exercise~1C12,
either $U_1\subseteq U_2$ or $U_2\subseteq U_1$. If the former occurs, then
$U_1\cup U_2=U_2$, which contains $U_3$. If the latter occurs, then
$U_1\cup U_2=U_1$. Since $U_3\subseteq U_2$, then $U_3\subseteq U_1$.
Therefore, one of the subspaces must contain the other two.

Now suppose no subspace is contained in another. Then there exist
$u\in U_2\setminus U_3$ and $v\in U_3\setminus U_2$. Since
$U_1\cup U_2\cup U_3$ is a subspace, $u+\alpha v\in U_1\cup U_2\cup U_3$ for
any nonzero $\alpha\in\F$. If $u+\alpha v\in U_2$, then
$\alpha v=(u+\alpha v)-u\in U_2$, which implies $v\in U_2$, a contradiction. If
$u+\alpha v\in U_3$, then $u=(u+\alpha v)-\alpha v\in U_3$, a contradiction.
Then $u+\alpha v\in U_1$ for all nonzero $\alpha\in\F$. Now choose nonzero
$\beta,\gamma\in\F$ with $\beta\ne\gamma$. Then
$u+\beta v,u+\gamma v\in U_1$, so
\[ (u+\beta v)-(u+\gamma v) = (\beta-\gamma)v\in U_1. \]
Since $\beta-\gamma\ne0$, we have $v\in U_1$. Similarly,
$(u+\beta v)-\beta v=u\in U_1$. Since $u\in U_2\setminus U_3$ and
$v\in U_3\setminus U_2$ were arbitrary, $U_2\setminus U_3\subseteq U_1$ and
$U_3\setminus U_2\subseteq U_1$.

Now let $w\in U_2\cap U_3$. Since $U_2\nsubseteq U_3$, there exists
$x\in U_2\setminus U_3\subseteq U_1$. Then $w+x\in U_2$. If $w+x\in U_3$, then
$x=(w+x)-w\in U_3$, a contradiction. Hence,
$w+x\in U_2\setminus U_3\subseteq U_1$, so $w=(w+x)-x\in U_1$. Therefore,
$U_2\cap U_3\subseteq U_1$, which gives
$U_2=(U_2\setminus U_3)\cup(U_2\cap U_3)\subseteq U_1$ and similarly
$U_3\subseteq U_1$, contradicting our assumption. Hence, one of the subspaces
must contain the other two.

\emph{Alternate proof.} \emph{Coolempire93.} Conversely, suppose that
$U_1\cup U_2\cup U_3$ is a subspace of $V$ but none of the subspaces contains
the other two. Let $W_1=U_2\cup U_3$ and $W_2=U_1\cup U_3$. By the assumption,
we cannot have $W_1\subseteq U_1$ or $W_2\subseteq U_2$, so there exist
$x\in U_1\setminus W_1$ and $y\in U_2\setminus W_2$. Now, since
$U_1\cup U_2\cup U_3$ is a vector space, we have
$x+\alpha y\in U_1\cup U_2\cup U_3$ for any nonzero $\alpha\in\F$. But if
$x+\alpha y\in U_1$, then $(x+\alpha y)-x=\alpha y\in U_1$, and as $\alpha$ is
nonzero, this means that $y\in U_1$, a contradiction. Similarly, if
$x+\alpha y\in U_2$, then $(x+\alpha y)-\alpha y=x\in U_2$, a contradiction.
Therefore, we must have $x+\alpha y\in U_3$ for all nonzero $\alpha$ to satisfy
the closure of the vector space.

As such, $x+y\in U_3$ and $x-y\in U_3$; that is, $(x+y)+(x-y)=2x\in U_3$ and
$(x+y)-(x-y)=2y\in U_3$, contradicting the original definitions of $x$ and
$y^\dagger$. Therefore, one of the subspaces must contain the other two.

$^\dagger$\,\emph{Final remark}. This is why the field must have more than two
elements. Since part of the field axioms is that $0\ne1$, if the field only has
two elements, it must be that $1=-1$ (consider addition mod $2$ as the addition
operation), so $2x\in U_3$ would not mean that $x\in U_3$, as $2x=0x=0$.
\end{solution}

\begin{exercise}
Suppose
\[ U = \{(x,-x,2x)\in\F^3 : x\in\F\} \quad\text{and}\quad
   W = \{(x,x,2x)\in\F^3 : x\in\F\}. \]
Describe $U+W$ using symbols, and also give a description of $U+W$ that uses no
symbols.
\end{exercise}
\begin{solution}
Let $u=(x,-x,2x)\in U$ and $w=(y,y,2y)\in W$. Then
\[ u+w = (x+y,-x+y,2x+2y) = \bigl(x+y,y-x,2(x+y)\bigr). \]
Observe that $x+y$ can take any value in $\F$, and $y-x$ can also take any value
in $\F$. Hence, we can write
\[ U+W = \{(a,b,2a)\in\F^3 \mid a,b\in\F\}. \]
In words, $U+W$ is the set of all vectors in $\F^3$ whose third component is
twice the first component, and whose second component can be any value in
$\F$.
\end{solution}

\begin{exercise}
Suppose $U$ is a subspace of $V$. What is $U+U$?
\end{exercise}
\begin{solution}
Let $u_1,u_2\in U$. Then $u_1+u_2\in U$ since $U$ is closed under addition.
Hence, $U+U\subseteq U$. Conversely, let $u\in U$. Then $u=u+0$, where $0$ is
the additive identity in $U$. Since $0\in U$, we have that $u\in U+U$. Hence,
$U\subseteq U+U$. Therefore, $U+U=U$.
\end{solution}

\begin{exercise}
Is the operation of addition on the subspaces of $V$ commutative? In other
words, if $U$ and $W$ are subspaces of $V$, is $U+W=W+U$?
\end{exercise}
\begin{solution}
Yes, the operation of addition on subspaces of $V$ is commutative. Let $u\in U$
and $w\in W$. Then $u+w\in U+W$. Because addition in $V$ is commutative, we
have that $u+w=w+u$, where $w\in W$ and $u\in U$. Hence, $w+u\in W+U$.
Therefore, every element of $U+W$ is also an element of $W+U$, which implies
that $U+W\subseteq W+U$. By a similar argument, we can show that
$W+U\subseteq U+W$. Hence, $U+W=W+U$.
\end{solution}

\begin{exercise}
Is the operation of addition on the subspaces of $V$ associative? In other
words, if $V_1,V_2,V_3$ are subspaces of $V$, is
\[ (V_1+V_2)+V_3 = V_1+(V_2+V_3)? \]
\end{exercise}
\begin{solution}
Yes, the operation of addition on subspaces of $V$ is associative. Let
$u_1\in V_1$, $u_2\in V_2$, and $u_3\in V_3$. Then
$(u_1+u_2)+u_3\in(V_1+V_2)+V_3$. Because addition in $V$ is associative, we
have that $(u_1+u_2)+u_3=u_1+(u_2+u_3)$, where $u_1\in V_1$, $u_2\in V_2$, and
$u_3\in V_3$. Hence, $u_1+(u_2+u_3)\in V_1+(V_2+V_3)$. Therefore, every element
of $(V_1+V_2)+V_3$ is also an element of $V_1+(V_2+V_3)$, which implies that
$(V_1+V_2)+V_3\subseteq V_1+(V_2+V_3)$. By a similar argument, we can show that
$V_1+(V_2+V_3)\subseteq(V_1+V_2)+V_3$. Hence, $(V_1+V_2)+V_3=V_1+(V_2+V_3)$.
\end{solution}

\begin{exercise}
Does the operation of addition on the subspaces of $V$ have an additive
identity? Which subspaces have additive inverses?
\end{exercise}
\begin{solution}
Yes, the operation of addition on the subspaces of $V$ has an additive
identity. The additive identity is the zero subspace $\{0\}$, which contains
only the zero vector of $V$. For any subspace $U$ of $V$, we have that
$U+\{0\}=U$.

If a subspace $U$ of $V$ had an additive inverse $W$, then we would have that
$U+W=\{0\}$. However, $U\subseteq U+W$ since $u=u+0$ for any $u\in U$. Hence,
$U+W=\{0\}$ implies that $U=\{0\}$. Therefore, the only subspace of $V$ that has
an additive inverse is the zero subspace $\{0\}$.
\end{solution}

\begin{exercise}
Prove or give a counterexample: If $V_1,V_2,U$ are subspaces of $V$ such that
\[ V_1+U = V_2+U, \]
then $V_1=V_2$.
\end{exercise}
\begin{solution}
The statement is false. Consider the vector space $\R^2$ and let
\[ V_1 = \{(x,0)\in\R^2 \mid x\in\R\}, \quad
   V_2 = \{(0,y)\in\R^2 \mid y\in\R\}, \quad
   U = \{(x,x)\in\R^2 \mid x\in\R\}. \]
Then $V_1+U=\R^2$ and $V_2+U=\R^2$, but $V_1\ne V_2$. Therefore, the statement
is false.
\end{solution}

\begin{exercise}
Suppose
\[ U = \{(x,x,y,y)\in\F^4 : x,y\in\F\}. \]
Find a subspace $W$ of $\F^4$ such that $\F^4=U\oplus W$.
\end{exercise}
\begin{solution}
Let
\[ W = \{(0,z,0,w)\in\F^4 \mid z,w\in\F\}. \]
If $(x,x,y,y)\in U$ and $(0,z,0,w)\in W$, then $(x,x,y,y)=(0,z,0,w)$ if and
only if $x=0$, $y=0$, $z=0$, and $w=0$. Hence, $U\cap W=\{0\}$. Moreover, for
any $(a,b,c,d)\in\F^4$, we can write
\[ (a,b,c,d) = (x,x,y,y)+(0,z,0,w), \]
where
\[ x=a,\quad y=c,\quad z=b-x=b-a,\quad w=d-y=d-c. \]
Hence, every vector in $\F^4$ can be written as a sum of a vector in $U$ and a
vector in $W$. Therefore, $\F^4=U\oplus W$.
\end{solution}

\begin{exercise}
Suppose
\[ U = \{(x,y,x+y,x-y,2x)\in\F^5 : x,y\in\F\}. \]
Find a subspace $W$ of $\F^5$ such that $\F^5=U\oplus W$.
\end{exercise}
\begin{solution}
Let
\[ W = \{(0,0,z,w,v)\in\F^5 \mid z,w,v\in\F\}. \]
If $(x,y,x+y,x-y,2x)\in U$ and $(0,0,z,w,v)\in W$, then
$(x,y,x+y,x-y,2x)=(0,0,z,w,v)$ if and only if $x=0$, $y=0$, $z=0$, $w=0$, and
$v=0$. Hence, $U\cap W=\{0\}$. Moreover, for any $(a,b,c,d,e)\in\F^5$, we can
write
\[ (a,b,c,d,e) = (x,y,x+y,x-y,2x)+(0,0,z,w,v), \]
which implies
\[ \begin{cases} a=x+0\\ b=y+0\\ c=(x+y)+z\\ d=(x-y)+w\\ e=2x+v \end{cases}
   \implies
   \begin{cases} x=a\\ y=b\\ z=c-(x+y)=c-(a+b)\\ w=d-(x-y)=d-(a-b)\\
     v=e-2x=e-2a \end{cases} \]
Hence, every vector in $\F^5$ can be written as a sum of a vector in $U$ and a
vector in $W$. Therefore, $\F^5=U\oplus W$.
\end{solution}

\begin{exercise}
Suppose
\[ U = \{(x,y,x+y,x-y,2x)\in\F^5 : x,y\in\F\}. \]
Find three subspaces $W_1,W_2,W_3$ of $\F^5$, none of which equals $\{0\}$, such
that $\F^5=U\oplus W_1\oplus W_2\oplus W_3$.
\end{exercise}
\begin{solution}
\solnote{Manual Remark}{The derivation of $z,w,e$ is from Exercise~1C21.}
Observe that $U$ can be expressed as
\[ U = \{(x,0,x,x,2x)\in\F^5 \mid x\in\F\}+\{(0,y,y,-y,0)\in\F^5 \mid y\in\F\}. \]
Define the subspaces
\begin{align*}
W_1 &= \{(0,0,z,0,0)\in\F^5 \mid z\in\F\},\\
W_2 &= \{(0,0,0,w,0)\in\F^5 \mid w\in\F\},\\
W_3 &= \{(0,0,0,0,v)\in\F^5 \mid v\in\F\}.
\end{align*}
Consider $(a,b,c,d,e)\in\F^5$. We can write
\begin{multline*}
(a,b,c,d,e) = (x,0,x,x,2x)+(0,y,y,-y,0)\\
  +(0,0,z,0,0)+(0,0,0,w,0)+(0,0,0,0,v),
\end{multline*}
where
\begin{align*}
x &= a,\\
y &= b,\\
z &= c-a-b,\\
w &= d-a+b,\\
v &= e-2a.
\end{align*}
Hence, every vector in $\F^5$ can be written as a sum of a vector in $U$ and
vectors in $W_1,W_2,W_3$. Moreover, if $a=b=c=d=e=0$, then $x=y=z=w=v=0$ so
that the only way to write the zero vector as a sum of vectors in
$U,W_1,W_2,W_3$ is to take all of them to be the zero vector. Therefore,
$\F^5=U\oplus W_1\oplus W_2\oplus W_3$.
\end{solution}

\begin{exercise}
Prove or give a counterexample: If $V_1,V_2,U$ are subspaces of $V$ such that
\[ V=V_1\oplus U \quad\text{and}\quad V=V_2\oplus U, \]
then $V_1=V_2$.
\exnote{Hint: When trying to discover whether a conjecture in linear algebra is
true or false, it is often useful to start by experimenting in $\F^2$.}
\end{exercise}
\begin{solution}
See Exercise~1C19 for a counterexample. For $V_1\oplus U$, observe that
$V_1\cap U=\{0\}$. Moreover, for any $(a,b)\in\R^2$, we can write
\[ (a,b) = (a-b,0)+(b,b)\in V_1\oplus U. \]
The proof for $V_2\oplus U$ is similar. Hence, $V=V_1\oplus U=V_2\oplus U$, but
$V_1\ne V_2$. Therefore, the statement is false.
\end{solution}

\begin{exercise}
A function $f\colon\R\to\R$ is called \emph{even} if
\[ f(-x) = f(x) \]
for all $x\in\R$. A function $f\colon\R\to\R$ is called \emph{odd} if
\[ f(-x) = -f(x) \]
for all $x\in\R$. Let $V_{\mathrm{e}}$ denote the set of real-valued even
functions on $\R$ and let $V_{\mathrm{o}}$ denote the set of real-valued odd
functions on $\R$. Show that $\R^{\R}=V_{\mathrm{e}}\oplus V_{\mathrm{o}}$.
\end{exercise}
\begin{solution}
Let $f\in\R^{\R}$. Define
\[ f_e(x) = \frac{f(x)+f(-x)}{2} \quad\text{and}\quad
   f_o(x) = \frac{f(x)-f(-x)}{2}. \]
Then $f_e$ is even since
\[ f_e(-x) = \frac{f(-x)+f(x)}{2} = f_e(x), \]
and $f_o$ is odd since
\[ f_o(-x) = \frac{f(-x)-f(x)}{2} = -\frac{f(x)-f(-x)}{2} = -f_o(x). \]
Moreover, we have that
\[ f(x) = f_e(x)+f_o(x). \]
Hence, every function in $\R^{\R}$ can be written as a sum of an even function
and an odd function.

Now, suppose $g\in V_e\cap V_o$. Then $g$ is both even and odd. Then for any
$x\in\R$, we have
\[ g(x) = g\bigl(-(-x)\bigr) = -g(-x) = -g(x), \]
which implies that $2g(x)=0$, or $g(x)=0$. Hence, the only function in the
intersection of $V_e$ and $V_o$ is the zero function. Therefore,
$\R^{\R}=V_e\oplus V_o$.
\end{solution}
